
%======================================================================
\section{v8: Ward combinations, the 34- and 18-dimensional maps, and \texorpdfstring{\(\phi\phi\)}{phi phi}}
\label{sec:v8map}
%======================================================================

\subsection{Constant-profile support of the Ward combinations}

Let \(K\) be the set of constant-profile coordinates, meaning jets
with \(m=0\). Let \(\pi_K\) be the coordinate projection onto
\(K\), and let \(\mathrm{Ann}_K\) be the annihilator of \(\pi_K Z\)
inside the span of the \(K\)-supported functionals. In every exact
run we found \(\abs K=12\), with HH~6, CB~2, CK~3 and PP~1. The
ST functionals supported on \(K\) span a space of dimension
\(6\). The \(m=0\) rows alone already span \(\mathrm{Ann}_K\).
Further, \(\rank\pi_KZ=6\) and \(\rank\pi_KB=5\). Also
\(\dim(Z\cap\ker\pi_K)=2\), except at \((7071/10^4,\Fr)\), where it
is \(3\).
The \(K\)-supported Ward combinations split by block support as
\[
 \underbrace{2}_{hh\text{ only}}+\underbrace{2}_{hh+\mathrm{CK}}
 +\underbrace{2}_{\mathrm{CB}+\mathrm{CK}}=6.
\]
No combination is supported on CB, CK or PP alone.
This decomposition is identical at every computed angle and
branch.

\subsection{The \texorpdfstring{\(hh\)}{hh} dictionary and the two exact Ward checks}

The engine's constant-profile \(hh\) coordinates are related to
the Wick-rotated longitudinal block by
\[
 x_{00}=-h_{\tau\tau},\quad x_{03}=\pm\ii\,h_{\tau z},\quad
 x_{33}=h_{zz},\quad f=\tfrac12(h_{11}+h_{22}),\qquad
 q_0=\pm\ii q_\tau,\ q_3=q_z.
\]
Both sign choices give a solvable linear system, and the two
give the same map. With
\(\Gamma_{\rm sym}=4\Gamma\), the six engine coordinates
\(y\) are
\[
 y=\bigl(-b_1,\,-2b_3,\,-b_2,\,b_6,\,b_7,\,-\tfrac14 b_8\bigr)
\]
in units of \(\ka^2/\pi\), where \(b_j\) are the coefficients of the
structures \eqref{eq:structures}
(from \(E=MS\) with \(4b=M^{\mathsf T}y\), recorded in
\loc{hh_map_*.json}).
The two \(hh\)-only Ward combinations are
\[
 W_1=b_1+2b_3,\qquad W_2=b_2 .
\]

\begin{proposition}[Exact \(hh\) Ward checks; executed]
\label{prop:v8hhward}
In exact arithmetic, at \(\alpha=3/10\) and at
\(\alpha=7071/10000\) and on both branches, the declared \(hh\) catalogue sum
satisfies \(W_1=W_2=0\).
The individual graphs do not. At \(\alpha=3/10\), for example,
\[
 (W_1,W_2)^{\rm D01}=\Bigl(-\tfrac{91}{5760},\ \tfrac{91}{2880}\Bigr)
 =-(W_1,W_2)^{\rm D03},
\]
and D02, D11 and D12 give zero.
\end{proposition}

\noindent
So the check is not trivially satisfied: it tests a cancellation
between D01 and D03. Like every catalogue statement, it is
conditional on Assumption~\ref{ass:reduction}.

\subsection{The remaining Ward combinations}

The two \(hh\)+CK and two CB+CK combinations contain the v7
ST row and the projected ghost equation restricted to \(m=0\)
(Theorems~\ref{thm:ghost} and \ref{thm:ST}). This identification
is \CONDITIONAL, because v8 builds no CK/CB coordinate dictionary.
The \(26-6=20\) combinations that involve \(m>0\) jets have not
been evaluated. They are \OPEN, not zero.

\subsection{The 18-dimensional engine map}

The HH sector of the inventory has rank 18 on the 34-candidate
branches, and rank 21 at \((7071/10^4,\Fr)\). The dictionary above
fixes the six constant-profile coordinates exactly, on all four
computed branches. The other twelve involve transverse
derivatives or non-axisymmetric jets, and their convention map is
\OPEN. The map is therefore \PARTIAL: 6 of 18 coordinates.

\subsection{The six exact values}

Write the non-PP part of a closed class as \(\sum_{i=1}^6e_iB_i\),
using the exact basis of \(B\) stored in
\loc{apex_cert_*_vectors.json}.
The \(hh\) constant-profile data give
\(\rank(\pi_{K\cap hh}B)=4\), and the resulting system is consistent:
\begin{center}
\small
\begin{tabular}{@{}lllll@{}}
\toprule
\((\alpha,\text{branch})\) & \(e_1\) & \(e_3\) & \(e_2+2e_5\) & \(e_4+e_5\)\\
\midrule
\((3/10,\CR)\) & \(7/40\) & \(0\) & \(7/120\) & \(-131/1920\)\\
\((3/10,\Fr)\) & \(7/40\) & \(0\) & \(7/120\) & \(63/640\)\\
\((7071/10^4,\CR)\) & \(2929/40000\) & \(0\) &
\(-32842877/282840000\) & \(-78439319/4525440000\)\\
%F7071ROW
\bottomrule
\end{tabular}
\end{center}
At both angles, \(e_1=(1-\alpha)/4\). This is observed at two
points, not proved for general \(\alpha\).
The CR/F difference appears only in \(e_4+e_5\), where it equals
\(\nu/8\) exactly: \(-1/6\) at \(\alpha=3/10\) and \(2071/28284\) at
\(\alpha=7071/10^4\), with \(\nu=4142/7071\).
This matches the \(hh\)-sum difference
\(b_8^{\CR}-b_8^{\Fr}=-\nu/4\), which is consistent with
\eqref{eq:branchdifference}.
The coefficients \(e_5\) and \(e_6\) are not determined. They need
CK/CB constant data, which in turn need a dictionary, and
\(m>0\) data, because \(\dim(Z\cap\ker\pi_K)=2\) (and \(3\) at
\((7071/10^4,\Fr)\)).
The status is \PARTIAL, and only in the stated basis.

\subsection{The \texorpdfstring{\(\phi\phi\)}{phi phi} values}

The first \(\phi\phi\) value is the catalogue PP coefficient
\(-c_2\ka^2P^2/(4\pi)\) on each branch. Its branch difference is
\(-\nu\ka^2P^2/(4\pi)\), and the engine coordinate is
\(y_{PP}=\ka^2c_2/(4\pi)\).
For the second value, the D05a part is \(0\) in the reduced model
of Section~\ref{sec:missing}.
The D04 part is still \OPEN, for the following reasons.
\begin{itemize}
\item At constant profile, the replay of D04 has a single key
\((0,0)\). It has no hidden mixed pair and no evanescent
coefficient, and its UNIT value is \(-\ka^2P^2/(4\pi)\).
\item The quadratic-profile part needs the \(r^2\)-weighted
mixed-channel trace. v8 proves only the ingredients for it
(Lemma~\ref{lem:v8duhamel}); the mode sum and the \(r^2\)
assembly were not done.
\end{itemize}
